% Istantanea a t = 0 dell'onda y = 0,030 sin(pi x - 20 pi t), in unità del SI
% (quesito ud4-onda-lettura-lambda): con k = pi rad/m la sinusoide si ripete
% ogni 2 m, la distanza fra due creste consecutive.
\begin{tikzpicture}[>=Stealth, font=\small, x=1.8cm, y=0.32cm]
  \draw[->] (0,0) -- (4.6,0) node[right] {$x$ (\si{\metre})};
  \draw[->] (0,-4) -- (0,5.6) node[above] {$y$ (\si{\metre})};
  \foreach \x in {1,2,3,4} { \draw (\x,0.25) -- (\x,-0.25); }
  \foreach \x in {1,3} { \node[anchor=north east] at (\x,-0.1) {$\x$}; }
  \foreach \x in {2,4} { \node[anchor=south east] at (\x,0.1) {$\x$}; }
  \draw (0.04,3) -- (-0.04,3) node[left] {$0{,}030$};
  \draw[formula, very thick, domain=0:4.4, samples=120] plot (\x, {3*sin(\x*180)});
  % creste e lunghezza d'onda
  \draw[dashed, richiamo] (0.5,3) -- (0.5,4.4);
  \draw[dashed, richiamo] (2.5,3) -- (2.5,4.4);
  \draw[|<->|] (0.5,4.4) -- (2.5,4.4) node[midway, above] {$\lambda = 2\pi/k = \SI{2}{\metre}$};
\end{tikzpicture}
